\DocumentMetadata{
  lang = pt-BR,
  pdfstandard = A-2b,
  pdfversion = 1.7
}

\documentclass{abntexto-ufc}

\ufcsetup{
  type = scientific-article,
  cover = false,
  catalog-card = false,
  coat-of-arms = false,
  glossary = none,
  index = none,
  author = {ARTICLEADVISORYAUTHOR},
  title = {Article Advisory Recommendation Evidence},
  submission-date = {1 de setembro de 2026},
  approval-date = {5 de setembro de 2026},
  article-author-note = {ARTICLEADVISORYNOTE Universidade Federal do Ceará. article@example.org}
}

\begin{document}

\ufcPrintArticleFrontMatter{%
  ARTICLEADVISORYSUMMARYOK apresenta uma síntese controlada para demonstrar o intervalo recomendado sem transformar orientação editorial em requisito de compilação. O cenário descreve um artigo acadêmico que organiza objetivo, método, resultados e conclusão em um único parágrafo, preservando linguagem clara, sequência lógica e vocabulário compatível com a área de normalização de trabalhos científicos. A análise considera que recomendações institucionais orientam a apresentação, mas não substituem requisitos técnicos expressamente classificados como obrigatórios no contrato normativo. Por isso, esta amostra mantém quantidade de palavras deliberadamente situada dentro do intervalo indicado para artigos e utiliza três palavras-chave, permitindo observar a configuração recomendada de forma positiva. O texto também registra que a classe deve continuar aceitando documentos válidos quando uma recomendação não for atendida, desde que nenhum requisito obrigatório seja violado. Essa distinção evita que preferências editoriais se tornem erros fatais e preserva a precedência de instruções específicas de periódicos quando o autor estiver preparando uma submissão externa. O cenário serve apenas como evidência controlada da modalidade recomendada.
  \ufcSummaryKeywords{normalização; artigo científico; evidência.}%
}

\end{document}
